\section*{Zielsetzung}
Die letzte Aufgabe für dieses Semester in STEP bestand darin eine dynamische Webanwendung zu erstellen, wo die von \textit{rhodes} empfangenen Daten angezeigt werden. Zudem sollen die Daten der jeweiligen Schiffe eine eigene Seite bekommen, welche auf der Hauptseite mit der MMSI des Schiffs verlinkt ist. Eine Sicherung der projekt-relevanten Skripte sowie die Erstellung eines Posters mit \LaTeX waren ebenfalls Teile des Projekts.
\\
Durch das erlernte Wissen aus STEP konnte das Projekt erfolgreich abgeschlossen werden. Von \textit{cat, mv, cd, ls} zu \textit{VAAK, Schleifen} und der \textit{Pipe} wurde nahezu alles genutzt, um das Projekt zu realisieren.
\\
\textbf{Ida Rhodes} (*1900 †1986) war eine Computerpionierin, welche als eine der ersten das Potenzial der Informationstechnik gesehen hat. In ihrem Vortrag '[t]he human computer's dream of the future' erzählte sie von ihrer Überzeugung der Nutzung dieser im Alltag~\autocite{Rhodes}. Ihr zu Ehren wurde der automatische Idetifikationssystem (kurz AIS) -Empfänger, welcher an der Hochschule genutzt wird, mit dem Namen \textit{rhodes} versehen.
